\documentclass[../main.tex]{subfiles} % 使用 subfiles 文档类，指定主文档
\graphicspath{{\subfix{../image/}}} % 指定图片目录，后续可以直接使用图片文件名。

% 例如:
% \begin{figure}[h]
% \centering
% \includegraphics{image-04.04}
% \label{fig:image-04.04}
% \caption{图片标题}
% \end{figure}

\begin{document}

\chapter{广义函数与Sobolev空间}

\begin{definition}
记多重指标$\alpha=(\alpha_1,\alpha_2,\cdots,\alpha_n)$，其中$\alpha_1,\alpha_2,\cdots,\alpha_n \geqslant 0$是整数,
\begin{align*}
|\alpha|=\sum\limits_{i=1}^{n}\alpha_i,\quad \alpha!=\alpha_1!\alpha_2!\cdots\alpha_n!,\\
x^\alpha=x_1^{\alpha_1}x_2^{\alpha_2}\cdots x_n^{\alpha_n}\quad (\forall x=(x_1,x_2,\cdots,x_n)\in\mathbb{R}^n),\\
\partial^\alpha=\partial_{x_1}^{\alpha_1}\partial_{x_2}^{\alpha_2}\cdots\partial_{x_n}^{\alpha_n},
\end{align*}
而且如果$\beta\leqslant\alpha$(系指$\beta_j\leqslant\alpha_j,j=1,2,\cdots,n$),则
\begin{align*}
\binom{\alpha}{\beta}=\frac{\alpha!}{\beta!(\alpha-\beta)!}=\binom{\alpha_1}{\beta_1}\binom{\alpha_2}{\beta_2}\cdots\binom{\alpha_n}{\beta_n}.
\end{align*}
\end{definition}

\subfile{Chapter-04/section-01.tex}

\subfile{Chapter-04/section-02.tex}

\subfile{Chapter-04/section-03.tex}

\subfile{Chapter-04/section-04.tex}

\subfile{Chapter-04/section-05.tex}

\subfile{Chapter-04/section-06.tex}

\subfile{Chapter-04/section-07.tex}

\subfile{Chapter-04/section-08.tex}

\subfile{Chapter-04/section-09.tex}

\subfile{Chapter-04/section-10.tex}

\end{document}